% Partie IR-CSG (méthode v2 : revue/methode_v2.md ; chiffres : sorties/architectures_ircsg.txt,
% sorties/controle_hauts_purs.txt ; tableaux générés par tables_ircsg.py)

L'impôt sur le revenu, la CSG et la CRDS prélèvent 220~Md€ sur les revenus soumis au barème : salaires,
pensions, bénéfices des indépendants, loyers, allocations chômage. Cette somme se décompose en 80~Md€ d'IR,
132~Md€ de CSG et 7~Md€ de CRDS. L'IR des revenus du capital au prélèvement forfaitaire n'y est pas compté :
ces revenus gardent leur régime (section~\ref{sec:pfu}). Les trois impôts se superposent sans se coordonner,
et la charge varie selon le profil :
\begin{itemize}\setlength{\itemsep}{0pt}
\item Pour une personne seule salariée, ils prélèvent 9,5~\% du salaire brut jusqu'à 23\,000~€ environ. Le
taux marginal saute alors à 21~\%, puis à 32~\% vers 40\,000~€ et à 40~\% au-delà de 110\,000~€
(figure~\ref{fig:baremes}).
\item Les retraités paient moins à revenu égal, du fait de l'abattement de 10~\% (5,3~Md€) et des taux
réduits de CSG. Ces taux réduits créent cependant des marches : quand le revenu fiscal franchit un seuil, le
nouveau taux s'applique à toute la pension, et 1\,000~€ de pension en plus peuvent coûter 620 à 770~€.
\item Le quotient familial des enfants réduit l'IR d'environ 13~Md€.
\end{itemize}

\subsection{Ce que dit la théorie}

Un taux marginal coûte par ce qu'il fait renoncer à gagner. Ce coût est d'autant plus élevé que l'euro en
question est déjà lourdement prélevé, et que les revenus concernés sont nombreux et capables de réagir.
\begin{itemize}\setlength{\itemsep}{0pt}
\item Une assiette large permet des taux plus bas là où ils coûtent. Les pensions déjà liquidées et les revenus
aujourd'hui exonérés par des niches (au moins 24~Md€) ne réagissent presque pas ; les taxer comme les salaires
permet, à recette égale, d'abaisser les taux qui pèsent sur le travail. C'est le mécanisme principal de la
fusion, et il fait des retraités les premiers contributeurs de la réforme.
\item L'élasticité des hauts revenus dépend des règles. Leur revenu déclaré réagit fortement au taux, mais
surtout par la forme, la date ou le lieu de la déclaration ; une assiette sans niche réduit ces possibilités.
Ce gain va à la recette, non à l'activité.
\item Les marches et les hauts taux du milieu coûtent plus que leur rendement ne le laisse croire. Entre 1 et
1,5~Smic, les cotisations non contributives et la sortie de la prime d'activité prélèvent déjà près de 60~\%
d'un euro de salaire en plus ; y ajouter un saut d'IR coûte beaucoup plus que ne le dit le seul taux de
l'impôt.
\item La décision de travailler dépend de l'impôt sur tout le premier salaire. Peu de salariés ajustent leurs
heures au taux marginal ; beaucoup plus choisissent de prendre un emploi, de reprendre après une interruption
ou d'ajouter un second revenu au foyer. Une franchise sur les premiers euros d'activité baisse l'impôt sur ce
premier salaire sans toucher au taux marginal des autres : c'est la logique du crédit d'impôt sur les revenus
d'activité introduit en Suède en 2007, dont les effets sur l'emploi restent d'ailleurs discutés.
\end{itemize}

\subsection{Ce qu'on peut chiffrer}

On mesure chaque architecture avec la méthode de l'encadré. Toutes partagent la même base :
\begin{itemize}\setlength{\itemsep}{0pt}
\item l'assiette est le revenu après cotisations, le même pour un salarié, un indépendant, un retraité ou un
chômeur, sans réduction, abattement ni taux réduit ;
\item l'impôt est individuel, et le même barème vaut pour tous ;
\item le taux marginal monte en ligne droite entre quelques seuils, jusqu'à 50~\% ;
\item la recette est toujours de 220~Md€.
\end{itemize}

\textbf{Le barème neutre.} Parmi ces barèmes, on retient d'abord celui qui s'écarte le moins de ce que chaque
foyer paie aujourd'hui : 2~\% au premier euro, 14~\% à 20\,000~€, 23~\% à 30\,000~€ et 50~\% à 250\,000~€ par
adulte (figure~\ref{fig:baremes}).
\begin{itemize}\setlength{\itemsep}{0pt}
\item \emph{Qui paie plus ou moins.} Aucun décile ne voit sa charge moyenne bouger de plus de 600~€
(tableau~\ref{tab:ircsg_dec}). Il reste qu'un foyer sur quatre paie plus de 500~€ de plus, à cause des niches
supprimées, de l'imposition individuelle et de l'alignement des retraités, qui paient 7~Md€ de plus.
\item \emph{Gain d'efficacité.} Le gain est de 4,0~Md€ par an, entre 1,9 et 9,2~Md€ selon les hypothèses.
\item \emph{D'où il vient.} Pour un peu moins de la moitié (1,8~Md€), de la participation : le taux d'entrée de
2~\% remplace les 9,7~\% de CSG et de CRDS prélevés dès le premier euro. Pour le reste, du dixième le plus
aisé : son taux marginal baisse, alors que son impôt moyen monte un peu avec la fin des niches.
\item \emph{Effet sur la recette.} La moindre optimisation des hauts revenus rapporterait en plus 1,5 à
4,5~Md€. Ce supplément est trop incertain pour financer quoi que ce soit, et il n'entre pas dans le gain.
\end{itemize}

\textbf{Ce qu'on peut obtenir au mieux.} Le barème simple qui minimise la perte sèche, franchise d'activité
comprise, montre ce que dirait la seule efficacité. Il sert ici d'étalon, non de proposition.
\begin{itemize}\setlength{\itemsep}{0pt}
\item Il gagne 7,8~Md€ (3,8 à 18,2).
\item Il est presque proportionnel : 16~\% jusqu'à 60\,000~€, et le taux de 50~\% n'est atteint qu'à 2~M€.
\item Le dernier décile y gagne 25~Md€, et la moitié des foyers perdent plus de 500~€ par an.
\end{itemize}

\figbloc{ircsg_p02_frontiere}{Gain d'efficacité et nombre de perdants : rendements décroissants}{fig:frontiere}

\textbf{Les rendements sont décroissants} (figure~\ref{fig:frontiere}).
\begin{itemize}\setlength{\itemsep}{0pt}
\item La fusion à répartition inchangée capte la moitié du gain atteignable.
\item Une franchise sur les revenus d'activité porte cette part à 59~\% avec 4\,000~€ par adulte, à 64~\% avec
8\,000~€. C'est le seul dispositif qui ajoute du gain, par la marge extensive.
\item Aller plus loin demande d'aplatir le barème au-dessus de 45\,000~€. Les 30 premiers pour cent du chemin
vers l'optimum portent la part de 52 à 69~\%, mais le dernier décile passe de +2,5 à $-$6~Md€. Le dernier tiers
du gain exige de lui rendre 20 à 25~Md€, et fait perdre plus de 500~€ à la moitié des foyers.
\end{itemize}
On obtient donc environ deux tiers du gain atteignable sans bouleverser la répartition par décile.

\begin{table}[!htbp]\centering\scriptsize\setlength{\tabcolsep}{3pt}
\caption{Architectures comparées, à recette égale (220~Md€)}\label{tab:ircsg_arch}
\resizebox{\linewidth}{!}{%
\begin{tabular}{@{}lrrrrrrr@{}}
\toprule
& Points & Gain d'efficacité & Part du & Perdants & Retraités & Couple, & 10~\% \\
& pour tous & Md€ par an & gain & $>$ 500~€ & Md€ & 2 enfants & du haut \\
& & (fourchette) & atteignable & & & € par foyer & Md€ \\
\midrule
Barème neutre & 0,0 & 4,0 (1,9 à 9,2) & 52~\% & 25~\% & +7,0 & +570 & +2,5 \\
\multicolumn{8}{@{}l}{\emph{Un dispositif ajouté au barème neutre}} \\
\quad Franchise d'activité 4~000~€ & +1,8 & 4,6 (2,1 à 10,2) & 59~\% & 30~\% & +14,0 & +217 & +3,6 \\
\quad Franchise d'activité 8~000~€ & +3,6 & 5,0 (2,2 à 11,0) & 64~\% & 33~\% & +21,4 & $-$172 & +4,7 \\
\quad Franchise de 8~000~€ sur tous les revenus & +5,8 & 3,3 (1,4 à 7,4) & 43~\% & 21~\% & +5,4 & +775 & +9,6 \\
\quad Franchise de 6~000~€ par enfant & +1,5 & 3,6 (1,7 à 8,4) & 47~\% & 27~\% & +13,0 & $-$983 & +0,6 \\
\quad Chèque de 742~€ par enfant & +0,9 & 3,1 (1,5 à 7,3) & 40~\% & 26~\% & +10,7 & $-$361 & +3,1 \\
\quad Quotient conjugal et familial & +2,0 & 3,9 (1,9 à 9,0) & 50~\% & 28~\% & +12,1 & $-$965 & $-$3,0 \\
\quad Abattement de 10~\% sur les pensions & +0,6 & 3,4 (1,6 à 8,0) & 44~\% & 24~\% & +0,9 & +923 & +2,7 \\
\quad Petites pensions exonérées & +0,4 & 3,7 (1,7 à 8,5) & 47~\% & 16~\% & +3,2 & +786 & +4,0 \\
\midrule
\textbf{Proposition}, franchise d'activité 4~000~€ & --- & 4,9 (2,3 à 11,1) & 63~\% & 29~\% & +15,2 & $-$1\,195 & $-$1,7 \\
Proposition, franchise d'activité 8~000~€ & --- & 5,2 (2,4 à 11,8) & 67~\% & 26~\% & +22,8 & $-$1\,629 & $-$1,5 \\
\midrule
Optimum d'efficacité, sans franchise & --- & 7,4 (3,7 à 17,4) & 95~\% & 58~\% & +13,7 & $-$519 & $-$25,7 \\
Gain atteignable (optimum, franchise 4~000~€) & --- & 7,8 (3,8 à 18,2) & 100~\% & 49~\% & +20,3 & $-$869 & $-$25,1 \\
\bottomrule
\end{tabular}}
\par\vspace{3pt}\parbox{\linewidth}{\raggedright\scriptsize\color{encre}
Gain d'efficacité : baisse de la perte sèche sur l'activité réelle par rapport à aujourd'hui, hypothèses centrales (basses et hautes entre parenthèses) ; la réaction du revenu déclaré n'y est pas comptée. Points pour tous : hausse des taux du barème neutre qui finance le dispositif. Proposition : franchise d'activité, 6~000~€ par enfant, pensions exonérées jusqu'à 12~500~€ puis de moins en moins jusqu'à 25~000~€ ; barème le plus efficace parmi ceux qui ne changent la charge moyenne d'aucun décile de plus de 500~€ et n'alourdissent pas celle du 1~\% le plus aisé. Retraités, couples, 10~\% du haut : variation de leur impôt (positif : ils paient plus).}
\end{table}


\textbf{Chaque dispositif a un prix} (tableau~\ref{tab:ircsg_arch}). Les dispositifs qui ne touchent pas à
l'activité, aide aux enfants ou ménagement des pensions, se financent par des taux un peu plus élevés pour
tous. Chacun réduit le gain de 0,1 à 0,9~Md€. Quatre arbitrages ressortent.
\begin{itemize}\setlength{\itemsep}{0pt}
\item \emph{Activité ou tous les revenus.} Une même franchise de 8\,000~€ étendue à tous les revenus coûte
deux points de plus et fait perdre un tiers du gain. Elle donne aux retraités ce que la franchise d'activité
réserve aux actifs, sans aucun effet sur l'emploi.
\item \emph{Ce que couvre la franchise d'activité.} Le modèle ne distingue pas, dans les revenus d'activité,
les salaires des loyers et des bénéfices. Il ne peut donc pas chiffrer une franchise réservée aux salaires.
Une franchise sur les loyers ne rend pourtant aucun travail plus payant.
\item \emph{Pensions.} Garder l'abattement de 10~\% ramène la hausse des retraités de 7 à 1~Md€, pour
0,6~point de taux pour tous. Exonérer les petites pensions la ramène à 3~Md€, pour 0,4~point.
\item \emph{Enfants.} Une franchise de 6\,000~€ par enfant est l'option la plus favorable aux familles :
1\,000~€ de moins en moyenne pour un couple avec deux enfants. Le chèque de 742~€ par enfant, qui reverse à
parts égales le coût actuel du quotient, l'est moins (360~€). Le quotient conjugal et familial profite surtout
aux couples aisés : le dixième le plus aisé y gagne 3~Md€.
\end{itemize}

\subsection{Ce que nous proposons}

\figbloc{ircsg_p01_baremes}{Taux moyen et taux marginal : aujourd'hui, barème neutre et proposition}{fig:baremes}

\textbf{La proposition} retient les trois dispositifs que la note défend, et fixe la répartition avant de
choisir le barème.
\begin{itemize}\setlength{\itemsep}{0pt}
\item \emph{Les dispositifs.} Une franchise de 4\,000~€ par adulte sur les revenus d'activité, pour l'emploi.
Une franchise de 6\,000~€ par enfant, déduite du revenu du parent qui gagne le plus, pour les familles.
L'exonération des pensions jusqu'à 12\,500~€, à peu près le minimum vieillesse, puis de moins en moins
jusqu'à 25\,000~€, pour que le minimum vieillesse ne soit pas imposé.
\item \emph{La contrainte de répartition.} La charge moyenne d'aucun décile ne doit bouger de plus de 500~€, et
celle du 1~\% le plus aisé ne doit pas monter.
\item \emph{Le barème.} Parmi ceux qui respectent la contrainte, on prend le plus efficace : 14~\% jusqu'à
20\,000~€ imposés, puis une montée en ligne droite jusqu'à 30~\% à 45\,000~€, puis une montée lente jusqu'à
50~\% à 1~M€.
\end{itemize}
Le gain d'efficacité est de 4,9~Md€ par an (2,3 à 11,1), soit 63~\% du gain atteignable ; 29~\% des foyers
perdent plus de 500~€.

\textbf{Qui paie.} Les actifs paient 15~Md€ de moins, les retraités 15~Md€ de plus. Le tableau~\ref{tab:ircsg_dec}
donne la variation par décile.
\begin{itemize}\setlength{\itemsep}{0pt}
\item \emph{Familles.} Un couple avec deux enfants paie 1\,200~€ de moins en moyenne (1\,360~€ à 50\,000~€ de
revenu), un parent isolé 230~€ de moins.
\item \emph{Salariés.} Un salarié au Smic paie à peu près ce qu'il paie aujourd'hui (+60~€) ; une personne
seule active, 320~€ de moins en moyenne.
\item \emph{Retraités.} Un retraité seul paie 600~€ de plus en moyenne, un couple de retraités 2\,300~€. Un
retraité seul à 17\,000~€ de pension, aujourd'hui exonéré, paierait 1\,260~€.
\end{itemize}
Entre 12\,500 et 25\,000~€ de pension, la sortie de l'exonération fait monter le taux marginal à 28~\%. C'est
sans effet sur l'activité, mais visible ; une sortie plus lente l'adoucirait, au prix d'un peu plus de taux pour
tous.

Avec une franchise de 8\,000~€, le gain monte à 5,2~Md€ (67~\%). Les retraités paient alors 23~Md€ de plus, et un
couple de retraités 3\,200~€. C'est l'arbitrage central de la réforme : l'emploi et les familles se financent
pour l'essentiel sur les pensions.

\begin{table}[!htbp]\centering\scriptsize\setlength{\tabcolsep}{3pt}
\caption{Variation de l'impôt par décile de revenu fiscal, toutes variantes (€ par foyer et par an)}\label{tab:ircsg_dec}
\resizebox{\linewidth}{!}{%
\begin{tabular}{@{}lrrrrrrrrrrrr@{}}
\toprule
& D1 & D2 & D3 & D4 & D5 & D6 & D7 & D8 & D9 & D10 & 1~\% (Md€) & Retraités (Md€) \\
\midrule
Réforme naïve (quotient, 3,99~\%) & $-$11 & $-$37 & +162 & +105 & $-$543 & $-$551 & $-$570 & +95 & $-$524 & +1\,841 & +2,6 & +8,2 \\
Réforme naïve, variante individuelle (1,64~\%) & $-$57 & $-$217 & $-$171 & $-$256 & $-$906 & $-$938 & $-$794 & +36 & +277 & +2\,953 & +2,8 & +2,7 \\
\textbf{Proposition} & $-$84 & $-$279 & +80 & +251 & $-$215 & $-$155 & $-$320 & +354 & $-$71 & +436 & 0,0 & +21,5 \\
Variante : franchise de 4\,000~€ & $-$85 & $-$149 & +147 & +225 & $-$286 & $-$129 & $-$237 & +304 & $-$158 & +367 & 0,0 & +14,3 \\
Variante : sans franchise & +2 & +19 & +243 & +318 & $-$225 & $-$202 & $-$414 & +113 & $-$259 & +399 & 0,0 & +6,9 \\
Quotient conjugal et familial & $-$84 & $-$285 & $-$18 & +119 & $-$276 & $-$137 & $-$149 & +468 & $-$99 & +460 & 0,0 & +18,9 \\
Abattement de 10~\% sur les pensions & $-$83 & $-$277 & +42 & +273 & $-$200 & $-$138 & $-$287 & +314 & $-$94 & +448 & 0,0 & +14,8 \\
Chèque par enfant (même coût) & $-$168 & $-$354 & +16 & +197 & $-$226 & $-$99 & $-$227 & +392 & +11 & +455 & 0,0 & +21,6 \\
Le 1~\% peut payer moins & $-$84 & $-$278 & +105 & +335 & $-$86 & $-$11 & $-$197 & +493 & +162 & $-$423 & $-$3,8 & +22,5 \\
Barème neutre (méthode v2) & 0 & +83 & +265 & +165 & $-$365 & $-$221 & $-$324 & +130 & $-$27 & +290 & $-$0,5 & +5,0 \\
Ancienne proposition (1~\% en baisse) & $-$85 & $-$145 & +190 & +342 & $-$150 & $-$12 & $-$121 & +450 & 0 & $-$453 & $-$2,5 & +13,9 \\
\bottomrule
\end{tabular}}
\par\vspace{3pt}\parbox{\linewidth}{\raggedright\scriptsize\color{encre}
Déciles de foyers fiscaux classés par revenu fiscal de référence (IRCOM, revenus 2024 ; montants de 2025). Moyennes de foyers représentatifs : elles cachent les écarts à l'intérieur de chaque décile (sous-partie~\ref{sec:ircsg}.4).}
\end{table}


\textbf{Les très hauts revenus.} Pour les 3\,800 foyers qui déclarent plus de 3~M€, le résultat dépend d'une
donnée que la DGFiP ne publie pas : la part de leurs revenus imposée au prélèvement forfaitaire.
\begin{itemize}\setlength{\itemsep}{0pt}
\item \emph{Si leurs revenus du barème paient aujourd'hui ce que dit le barème légal}, 44 à 49~\% : la
proposition baisse leur impôt de 6 à 13~points de ce revenu. Seuls les 700 foyers au-delà de 9~M€ paient
1~point de plus.
\item \emph{Dans la convention centrale du modèle}, qui leur prête un taux actuel de 32 à 37~\%, difficile à
concilier avec le barème : ils paieraient 2 à 16~points de plus.
\end{itemize}
Dans les deux cas, la comparaison omet la CEHR, qui disparaîtrait sur les revenus du barème et leur retirerait
3 à 4~points. Leurs revenus du capital restent au prélèvement forfaitaire, majoration comprise
(section~\ref{sec:pfu}). Caler cette part sur les données de la DGFiP par tranche est le premier travail à faire
avant de chiffrer définitivement.

\textbf{La transition.}
\begin{itemize}\setlength{\itemsep}{0pt}
\item La collecte ne change presque pas : employeurs, caisses de retraite et France Travail prélèvent déjà à
la source l'IR et la CSG sur les mêmes revenus. Le taux individualisé est la règle par défaut du prélèvement à
la source depuis le 1\textsuperscript{er}~septembre~2025 (loi de finances pour 2024, art.~19).
\item Une convergence linéaire sur quatre ans limiterait la hausse annuelle du retraité à 17\,000~€ à environ
320~€. Comme gagnants et perdants avancent au même rythme, la recette reste à peu près constante chaque année.
\end{itemize}

\textbf{Le risque constitutionnel.} Le Conseil constitutionnel a censuré :
\begin{itemize}\setlength{\itemsep}{0pt}
\item en 2000, une réduction de CSG qui ignorait les autres revenus du foyer et les charges de famille
(décision 2000-437~DC) ;
\item en 2012, une contribution assise sur chaque personne sans tenir compte du foyer (2012-662~DC).
\end{itemize}
Un impôt individuel avec une franchise réservée aux revenus d'activité réunit ces deux traits. La proposition
s'en distingue parce qu'elle remplace l'impôt de foyer au lieu de s'y ajouter et qu'elle tient compte des
enfants. L'argument reste à construire. La variante « quotient », qui rétablit l'imposition commune des
couples, est la solution de repli : elle coûte peu en efficacité (3,9~Md€ de gain au lieu de 4,0 pour le
barème neutre).
